\chapter{学习方法：从新手到专家的路径} \label{app:learning-methodology}

本书正文讨论的是“写什么代码”和“怎么调试”，但读者在阅读过程中问得最多的往往是另一类问题：“我应该先学什么？怎么学效率最高？为什么我看了书还是不会做项目？”本章不讨论具体技术，只讨论如何学习嵌入式开发这件事本身。作者结合自己从小学机器人兴趣班到工作的完整历程，提炼出一套可复用的学习方法论。

本章的核心观点只有四条：
\begin{enumerate}
	\item 学习是分阶段的，不同阶段需要不同的策略——不要用入门阶段的方法去啃专家阶段的内容；
	\item 阅读代码是比写代码更重要的能力——看不懂别人的代码，就写不出自己的代码；
	\item 调试是学习的核心环节——每一次崩溃都是学习机会；
	\item 用 AI 加速学习，而不是用 AI 替代思考。
\end{enumerate}

\section{学习的四个阶段}

作者将嵌入式学习分为四个阶段。每个阶段的目标、方法和常见误区都不同。

\subsection*{阶段一：模仿（0——6 个月）}

目标：能跑通别人的代码，知道“按这个步骤做就能看到结果”。

典型行为：照着教程一步步操作，把例程下载到开发板，看到 LED 闪烁或串口打印就感到兴奋。遇到报错时复制错误信息去搜索引擎，找到解决方案后照做，不一定理解为什么。

有效方法：
\begin{itemize}
	\item 找一份完整的、带注释的例程，从头到尾跑一遍，不要跳步；
	\item 每次修改代码只改一处，观察变化，建立“代码改动 → 现象变化”的因果链；
	\item 遇到报错时，先尝试自己读错误信息，再搜索——搜索本身也是学习；
	\item 用 Git 做版本管理，每跑通一个功能就 commit 一次，方便回退。
\end{itemize}

常见误区：
\begin{itemize}
	\item 只看不练：看十遍教程不如亲手敲一遍代码；
	\item 追求“理解每一行”再跑：在模仿阶段，先跑通再理解是更有效的顺序；
	\item 频繁切换开发板：这个阶段最重要的是“把一块板子玩透”，而不是“玩过很多块板子”。
\end{itemize}

\subsection*{阶段二：迁移（6——18 个月）}

目标：能把 A 项目中的代码搬到 B 项目中使用，能在现有代码基础上添加新功能。

典型行为：不再满足于跑通例程，开始尝试“把这个 LCD 驱动移植到我的项目里”、“把那个按键扫描算法改成我需要的模式”。开始理解模块化设计的好处——因为模块化程度高的代码更容易迁移。

有效方法：
\begin{itemize}
	\item 主动给自己出题：比如“把例程里的串口打印改成 LCD 显示”、“把按键轮询改成中断方式”。这些题目有明确答案但需要自己动手；
	\item 收集整理自己验证过的代码片段，形成个人代码库（对应正文第 1 章“构建自己的代码库”）；
	\item 对比不同项目的实现方式：同一个功能（如 FIFO 队列），不同作者可能有不同写法，对比能加深理解；
	\item 开始阅读芯片参考手册的“相关章节”，而不是只依赖教程。
\end{itemize}

常见误区：
\begin{itemize}
	\item 迷信“重写才算是自己的”：在别人代码基础上修改、适配、扩展，是工程开发的常态，不必每次从零开始；
	\item 只移植不总结：移植完一个模块后，花 10 分钟写几句笔记——“踩了什么坑、为什么在这里改了一行”，这 10 分钟的 ROI 极高；
	\item 害怕改别人的代码：在迁移阶段，要敢于修改例程代码，改崩了就用 Git 回退。
\end{itemize}

\subsection*{阶段三：建造（18——36 个月）}

目标：能独立完成一个完整项目，从需求分析到硬件选型到代码实现到调试交付。

典型行为：不再只是“用别人的框架写代码”，而是开始思考“这个项目需要什么样的框架”。开始关注软件架构、可维护性、测试策略。能在多个方案中做取舍——用轮子还是自研、用 RTOS 还是裸机、用状态机还是回调。

有效方法：
\begin{itemize}
	\item 完整地做一个项目——不是“跑通例程”，而是从零开始完成一个“能用”的东西。项目不必大，但必须完整（有需求、有设计、有实现、有测试、有交付）；
	\item 项目完成后写一份“复盘文档”：哪些做对了、哪些可以改进、如果重来会怎么设计。这份文档比代码本身更有价值；
	\item 读别人项目的源码，尤其是那些“你用过但没读过”的库（如 FreeRTOS 的部分源码、FatFS 的核心逻辑）；
	\item 参与技术社区的讨论，尝试解答别人的问题——这是检验自己理解深度的好方法。
\end{itemize}

常见误区：
\begin{itemize}
	\item 追求“新技术”而忽视“基本功”：在这个阶段，精通一个领域比浅尝十个领域更有价值；
	\item 过度设计：为“未来可能有的需求”写了很多抽象层，结果项目交付时这些抽象层从未被使用；
	\item 不写文档：这个阶段的项目通常较复杂，不写文档的话三个月后自己都看不懂。
\end{itemize}

\subsection*{阶段四：创造（36 个月以上）}

目标：能解决别人没解决过的问题，能设计新的方案、框架或工具。

典型行为：不再满足于“用现有方案解决问题”，而是思考“有没有更好的方法”。开始写工具来提升效率（如第~\ref{ch:script-engine}~章的解析器生成器、附录~\ref{app:author-projects}~的 MetroEditor）。能把一个领域的方法迁移到另一个领域。

有效方法：
\begin{itemize}
	\item 选择一个“痛点”深挖：比如“为什么每次调 PID 都要重编译？”“为什么串口调试总丢帧？”——解决一个真正的痛点，就创造了一个有价值的东西；
	\item 把经验抽象成方法论：能说出“在这个场景下应该用 A 方案而不是 B 方案，因为……”——不再只是“我知道怎么做”，而是“我知道为什么这么做”；
	\item 输出：写文章、写书、做开源项目。输出的过程会迫使你整理知识体系，发现自己的盲点。
\end{itemize}

常见误区：
\begin{itemize}
	\item 闭门造车：创造力不是凭空产生的，它需要大量输入（读论文、读别人的代码、参加技术交流）；
	\item 忽视工程约束：学院派的“最优方案”在工程中可能因为成本、交期、团队能力而不适用。真正的创造力是在约束条件下找到可行解。
\end{itemize}

\begin{table}[htbp]
	\centering
	\caption{学习四阶段对比}
	\label{tab:app-learning-stages}
	\begin{tabularx}{\textwidth}{p{2.2cm}p{2.8cm}p{2.8cm}p{2.8cm} X}
		\toprule
		& 模仿 & 迁移 & 建造 & 创造 \\
		\midrule
		核心任务 & 跑通例程 & 移植复用 & 完整项目 & 解决新问题 \\
		主要方法 & 照做、观察 & 对比、适配 & 设计、取舍 & 抽象、输出 \\
		依赖资源 & 教程、例程 & 代码库、手册 & 项目经验 & 领域知识、方法论 \\
		典型产出 & 能亮灯 & 能复用模块 & 能交付产品 & 能创造工具/框架 \\
		\bottomrule
	\end{tabularx}
\end{table}

\section{如何阅读技术文档}

初学者面对上千页的芯片参考手册时感到无从下手。作者的建议是：不要从头读到尾，而是带着问题去读。

\subsection{数据手册 vs 参考手册}

STM32 系列有两份核心文档：

\begin{itemize}
	\item 数据手册：描述芯片的“硬件特性”——引脚定义、电气参数、封装尺寸、温度范围。选型时主要看这份；
	\item 参考手册：描述芯片的“功能细节”——寄存器定义、外设工作模式、时序要求。编程时主要看这份。
\end{itemize}

初学者常犯的错误是：把参考手册当教科书从头读到尾。正确做法是：

\begin{enumerate}
	\item 先浏览目录，知道“这个芯片有哪些外设”；
	\item 需要用哪个外设时，再去读对应章节；
	\item 第一次读一个外设时，重点关注“功能概述”和“寄存器描述”，暂时跳过“时序图”和“应用场景”；
	\item 配合例程一起看：手册里描述了“这个寄存器控制什么”，例程里展示了“实际代码怎么写”。
\end{enumerate}

\subsection{如何阅读别人的代码}

阅读代码是比写代码更重要的能力。建议按以下顺序阅读一个陌生的代码库：

\begin{enumerate}
	\item 先读 README：了解这个项目是做什么的、怎么编译、怎么运行；
	\item 找入口：对于嵌入式项目，入口通常是 \texttt{main.c} 里的 \texttt{main()} 函数或启动文件中的复位向量；
	\item 追踪主流程：从入口开始，沿执行路径往下读，不要深入每个函数，先建立“这个程序大概做了什么”的心智模型；
	\item 读核心数据结构：找到项目中最重要的结构体（如任务控制块 TCB、菜单项结构体），理解它们的关系；
	\item 读关键函数：选择 3——5 个最核心的函数深入阅读，理解它们的输入、输出和副作用；
	\item 读测试用例：如果有测试代码，读测试能最快理解“这个函数应该怎么用”。
\end{enumerate}

读代码时不建议在代码里直接加注释（可能污染原仓库），建议单独建一个 \texttt{notes/} 目录，每个文件对应一个模块，记录：
\begin{itemize}
	\item 这个模块的职责是什么；
	\item 它的对外接口有哪些；
	\item 它依赖哪些其他模块；
	\item 你读的时候遇到了什么问题、是怎么解决的。
\end{itemize}

\section{如何有效记笔记}

关于“要不要记笔记”，作者的观点是：核心笔记要写，但笔记的目的不是“存档”，而是“激活思考”。

\subsection{两种笔记系统}

\begin{itemize}
	\item 收集型笔记：摘抄、收藏、剪藏。这类笔记的“消化率”很低——收藏了 ≠ 学会了。作者的经验是：收集型笔记只占 10\% 的精力，不必追求“全面”；
	\item 产出型笔记：用自己的话重新组织、总结、提炼。比如“读完一个外设章节后，用一页 A4 纸画出它的状态机”、“修复一个 Bug 后写一段文字描述根因和修复方案”。这类笔记的“消化率”高，应该占 90\% 的精力。
\end{itemize}

\subsection{为什么要写“问题和答案”笔记}

作者发现，最有效的笔记格式是“我遇到了什么问题 → 我是怎么解决的 → 为什么这个方案可行”。这种笔记有几个好处：

\begin{enumerate}
	\item 真实：问题来自实际调试，不是教科书上的理想案例；
	\item 可检索：未来遇到类似问题时，通过关键词能快速定位；
	\item 可复用：解决方案本身就是经验积累。
\end{enumerate}

示例（来自作者调试串口 FIFO 的笔记）：

\begin{quote}
	问题：串口接收中断中，从 FIFO 读数据时偶尔丢失字节。
	
	现象：高波特率（115200）下，连续发送 100 字节，接收端有时只收到 97 字节。
	
	根因：中断中同时修改了头指针和尾指针，主循环也在修改尾指针，存在竞争条件。
	
	解决：在中断和主循环中分别加入临界区保护（关中断），用 \texttt{\_\_disable\_irq()}/\texttt{\_\_enable\_irq()} 保护对 FIFO 头尾指针的修改。
	
	验证：连续发送 10000 字节，接收端全部正确。
\end{quote}

这类笔记积累到一定数量后，就成了个人的“调试手册”。正文第~\ref{ch:deadline-ai}~章提到的 \texttt{MEMORY.md} 就是这种思路的延伸——把个人经验固定下来，让 AI 也能了解你的习惯。

\section{如何调试——学习的最佳时机}

调试不是学习的“副产品”，而是学习的“核心环节”。每次程序崩溃、每次现象不符预期，都是深度理解系统运行机制的机会。

\subsection{调试的心态}

\begin{itemize}
	\item 把 Bug 当成线索：程序不会“无缘无故”出错。每一个异常都是“程序实际运行状态”与“你以为的运行状态”之间的偏差。这个偏差本身就是信息；
	\item 不要迷信“经验”：“我以前遇到过类似问题”是起点，不是终点。每次都要验证根因，而不是凭记忆“大概改一下”；
	\item 一次只改一处：调试时一次修改多个地方是效率最低的做法——改完后你不知道是哪处修改起了作用。一次只改一处，改完重新测试。
\end{itemize}

\subsection{调试的流程}

作者推荐的调试流程：

\begin{enumerate}
	\item 复现：确保 Bug 能稳定复现。“偶尔出现”的问题先尝试找到复现条件；
	\item 定位：缩小问题范围——是硬件还是软件？是中断还是主循环？是初始化阶段还是运行阶段？用二分法加日志逐步缩小范围；
	\item 理解：定位到出错的代码行后，不要急着改。先问“为什么这行会出错”——是前面有状态没设置？是某个变量被意外修改了？是时序不对？
	\item 修复：理解根因后，设计修复方案。评估修复方案的副作用——这个改动会影响其他功能吗？
	\item 验证：修复后重新跑测试，不仅跑复现场景，还要跑相关的边缘场景。确保没有引入新 Bug；
	\item 记录：把调试过程和根因写下来。这是最有价值的学习材料。
\end{enumerate}

\subsection{常用的调试工具}

\begin{table}[htbp]
	\centering
	\caption{调试工具与适用场景}
	\label{tab:app-debug-tools}
	\begin{tabularx}{\textwidth}{lXX}
		\toprule
		工具 & 适用场景 & 优缺点 \\
		\midrule
		串口打印（printf） & 查看变量值、状态变化、执行路径 & 简单通用，但干扰时序 \\
		硬件断点（JTAG/SWD） & 查看特定位置时的寄存器状态 & 精确，但需要调试器 \\
		逻辑分析仪 & 查看数字信号时序（I2C、SPI、UART） & 能看到物理层，但需要额外设备 \\
		示波器 & 查看模拟信号（电压、PWM、电源纹波） & 物理世界唯一窗口，需要示波器 \\
		计数器/定时器 & 测量代码执行时间 & 精确，无侵入 \\
		\texttt{assert} & 在运行时检查不变量 & 低成本快速检测，需确保不被优化 \\
		\bottomrule
	\end{tabularx}
\end{table}

串口打印是最常用的调试手段，但很多人用不好。几个进阶技巧：

\begin{enumerate}
	\item 带时间戳：每条日志加一个时间戳（如 \texttt{msTick}），可以判断事件发生的时序；
	\item 分级日志：TRACE、DEBUG、INFO、WARN、ERROR 分级，调试时只开需要的级别，避免信息过载；
	\item 用宏控制：\texttt{\#ifdef DEBUG\_ENABLE} 控制调试代码是否编译，避免发布版本中残留过多日志；
	\item 环形缓冲区 + 触发保存：把日志写入 RAM 中的环形缓冲区，只有在发生异常时才把缓冲区内容通过串口输出。这样能看到崩溃前的完整状态，不干扰实时性。
\end{enumerate}

\subsection{“塑料轮子”在学习中的意义}

正文第~\ref{ch:deadline-ai}~章和第~\ref{ch:about-author}~章反复提到“塑料轮子”的概念。从学习方法的角度看，“塑料轮子”还有另一层含义——它是验证理解的最佳工具。

当你觉得自己理解了某个原理，可以尝试“造一个塑料轮子”来验证：
\begin{itemize}
	\item 理解状态机了？试着写一个最简单的按键状态机（10 行代码），不要用任何库；
	\item 理解 FIFO 了？试着写一个环形缓冲区（20 行代码），自己实现入队和出队；
	\item 理解 PID 了？在 PC 上写一个仿真程序（50 行 Python），观察积分饱和的效果。
\end{itemize}

这些“塑料轮子”虽然简陋，但它们有一个共同特点：每一行代码都是你亲手写的，每一行都经过你的验证。造过之后，再用成熟的库或框架时，你不再“盲信”它的正确性——你知道它在背后大概做了什么，也知道它可能在哪里出问题。

\section{如何利用 AI 加速学习}

正文第~\ref{ch:deadline-ai}~章讨论了在交期压力下如何使用 AI 工具，这里补充学习方法层面的建议。

\subsection{AI 能加速什么}

\begin{itemize}
	\item 解释代码：把一段你看不懂的代码贴给 AI，让它逐行解释。比搜索引擎更直接；
	\item 生成示例：让 AI 生成“展示某个函数用法的简短示例代码”，比自己从头翻文档快；
	\item 解释错误信息：把编译错误或运行错误的输出贴给 AI，它能快速分析可能原因；
	\item 对比方案：问 AI “方案 A 和方案 B 各有什么优缺点”，比自己翻多份文档高效。
\end{itemize}

\subsection{AI 不能替代什么}

\begin{itemize}
	\item 动手调试：AI 可以分析日志、提出建议，但它不能帮你插示波器、不能帮你按复位键、不能帮你感受“电机转起来有没有异响”；
	\item 设计取舍：AI 可以列出多个方案的技术优缺点，但“在这个交期下、这个团队能力下、这个成本约束下，选哪个方案”需要你自己判断；
	\item 建立直觉：AI 可以解释“I2C 的时序要求”，但不能替代你在示波器上亲眼看到波形不对时的那种“感觉”——经验只能从实践中来。
\end{itemize}

\subsection{学习时如何用 AI（推荐做法）}

\begin{enumerate}
	\item 先自己思考，再问 AI：遇到问题先自己分析 5 分钟——哪怕结论是“我不确定”，这个思考过程本身有价值。不要一遇到问题就转向 AI；
	\item 让 AI 解释，不替 AI 相信：AI 给出的解释要自己验证。比如 AI 说“这里应该用 volatile”，你要自己去查 volatile 的含义，确认它说得对不对；
	\item 用 AI 做“加速器”，不做“拐杖”：AI 可以帮你节省查阅资料的时间，但它不能替你建立知识体系。最终对项目负责的是你，不是 AI。
\end{enumerate}

\section{学习心态}

\subsection{容忍“暂时不懂”}

嵌入式学习的一个特点是：知识之间有依赖关系。你可能需要先理解“中断”才能理解“串口接收”，需要先理解“指针”才能理解“FIFO 队列”。在学习初期，很多概念暂时看不懂是正常的。

作者的建议是：遇到暂时看不懂的概念，做一个标记，先跳过去，等后续有了更多上下文再回来看。不要因为一个概念卡住而停止整个学习进程。

\subsection{接受“不完美”的交付}

在学习和项目中，完美主义往往是效率的敌人。作者见过太多人卡在“要把这个例程完全理解了再进入下一章”，结果看了三章就放弃了。

更好的策略是：先跑通、再理解、再优化。跑通给你“反馈”——看到结果会增强信心；理解给你“能力”——知道为什么能跑；优化给你“经验”——知道怎么做得更好。

\subsection{输出是最好的学习方式}

作者坚持认为：只输入不输出，学习效果会打折扣。输出可以有多种形式：

\begin{itemize}
	\item 写学习笔记（哪怕只是给自己看）；
	\item 写一篇技术博客（在论坛、公众号或自己的博客）；
	\item 回答别人的问题（在技术社区或学习群）；
	\item 做一个开源的小项目（哪怕只是一个“能亮的灯”）。
\end{itemize}

输出的过程会迫使你“把模糊的理解变成清晰的表达”，这个过程本身就在加深理解。

\section{总结本章}

学习方法的核心不是“学得快”，而是“学得有效”。不同阶段用不同的策略——模仿阶段多动手、迁移阶段多对比、建造阶段做完整项目、创造阶段做抽象输出。阅读技术文档和代码要带着问题，调试是学习的最佳时机，AI 是加速器但不是拐杖。最重要的是，保持输出和分享的习惯——因为它不仅能让别人受益，更能让你自己受益。

作者在写这本书的过程中，最大的收获不是“把知识整理成了文字”，而是“为了写清楚，发现自己原来还有没想透的地方”。愿每一位读者都能找到适合自己的学习方法，在自己热爱的领域里，一步一个脚印地走下去。
